\section{Formal Falsification}
\label{ch:falsify}

The falsification proceeds in five steps. Ill-posedness decides that inference cannot do without context; this was already accomplished in Section~\ref{ch:formalize}. The remaining four steps answer four questions in turn: whether context strength can be governed uniformly by one global scalar, whether it can be appended outside the generative objective, what the open set demands of the source of the rule's values, and on which side of the decision axis the two canonical components fall.

\emph{Reading convention.} Each item below is proved in the \emph{Argument} paragraph printed with it, at the level of generality stated there, so that no named claim has to be reconstructed from context. Where the machine-checked form of an item differs from the prose --- stronger, weaker, or resting on different premises --- Section~\ref{ch:verify} records the difference explicitly, item by item.

A word on the nontriviality of this chain. Two steps of the chain are not contained in the definitions: Lemma 1's exclusion of the global-scalar form is an algebraic fact about optimality conditions, and Lemma 2's exclusion of the frozen extrinsic paradigm is likewise. The transmission from distributional openness to Mode A is indeed close to the definition level, and making this transmission explicit is exactly the content of the decision axis; but the two ends of the transmission, namely the form exclusion and the classification closure, are independent mathematical content, and the force of the falsification is carried by them.

The logical anatomy of the negation should be stated first. Criteria P1 to P3 are structural clauses at the task level, characterized item by item by the mask form and the inverse-problem setting; Criterion P4 is the sole premise of an empirical kind, and Section~\ref{sec:p4} states it in two components, a barrier condition asserting that the demand separates at every scale without a pre-given modulus, and a realization condition asserting that the separated instances carry positive probability under the test distribution. Of the three premises of Lemma 3, upper-tail realization is a theorem rather than a premise, and nonnegativity of the minimum risk is trivial under squared loss; neither carries empirical content. The force of Theorem 5 therefore equals exactly the force of P4: no weaker and no stronger. This equality is the paper's declaration of its own conditionality rather than a retreat: a universal negation over a task class must have its load-bearing premise either in the task structure, here P1 to P3, or in an empirical assertion about the task, here P4; No-Free-Lunch carries its universal law on the premise of uniform averaging, and this paper correspondingly carries its negation on a checkable criterion. What the decomposition adds is that P4 is not a restatement of the conclusion: the barrier condition mentions no rule at all, and it fails on a task that satisfies P1 to P3; the derivation from the two components, together with the bounded-description condition that a fixed parameter object satisfies, to the risk-gap floor of \eqref{eq:floor} is a proved inequality rather than a definitional unfolding. The quantified object of P4 is the class of all functions computable by parameters fixed at inference time, that is, every rule whose values pre-exist inference; the only regularity asked of that class is the bounded-description condition of Section~\ref{sec:p4}, which says that a finite description responds to the instance at a finite rate, and which by the implementer clause of Definition~3 sets aside only objects whose values are analytic invariants and therefore belong to Mode A in the first place.

\subsection{Lemma 1: A Single Global Scalar Cannot Realize a Heterogeneous Strength Profile}

The most economical approach is to govern context strength uniformly by one global scalar $\lambda$ --- in the indexed notation of Definition~2, the constant family $R_i \equiv \lambda$ --- which is exactly the form of classical regularization \citep{poggio1985}. The classical regularization tradition has long known the limitation of fixed strength and has developed parameter-selection methods such as the L-curve and generalized cross-validation \citep{hansen1992,golub1979}. Lemma 1 of this section algebraizes this limitation into a contradiction at the level of optimality conditions and further points out that any parameter selection completed offline still produces a frozen parameter table and belongs to Mode B. Let the data term be the per-position squared residual $\ell_i (x, d) = (x_i - d_i)^2$, and let the per-position coefficient of the consistency term be
\begin{equation}
\kappa_i (d) = \sum_{j \in I} (d_i - d_j),
\label{eq:kappa}
\end{equation}
which measures the tendency of position $i$ to deviate from its neighborhood; a large $\kappa_i (d)$ indicates that the position tends to be pulled toward the neighborhood. The objective with extrinsic smoothing is written as
\begin{equation}
\mathcal{E}(x, d, \lambda) = \sum_{i \in I} (x_i - d_i)^2 + \lambda \sum_{i \in I} \sum_{j \in I} (d_i - d_j)^2.
\label{eq:paradigmE}
\end{equation}
Differentiating with respect to $d_i$ and normalizing the constant factors, the first-order optimality condition is
\begin{equation}
(d_i - x_i) + \lambda \cdot \kappa_i (d) = 0.
\label{eq:foc}
\end{equation}
The meaning of \eqref{eq:foc} is: at every position, the weighted sum of the data residual and the consistency term must vanish. Now ask: can one and the same $\lambda$, that is a constant family, serve two positions whose required strengths differ?

\begin{lemma}
If the context rule is global, so that the family of Definition~2 is constant, $R_i = \lambda$ for every position $i \in I$, while the task demands different coupling strengths at boundary and homogeneous positions, then the scalar cannot satisfy the first-order optimality conditions at both ends: the heterogeneous demand breaks the cross-position proportionality, and no constant family realizes a heterogeneous strength profile.
\end{lemma}

\argument Lemma 1 has a strengthened form, covering arbitrary fitting points rather than only exact fits, whose premise is that the following proportionality fails to hold. Suppose the same scalar $\lambda$ satisfies the first-order conditions of position $i$ and position $j$ simultaneously at an arbitrary point $d$. From \eqref{eq:foc} one obtains $\lambda \cdot \kappa_i (d) = x_i - d_i$ and $\lambda \cdot \kappa_j (d) = x_j - d_j$ respectively. Cross-multiplying the two equations to eliminate $\lambda$ gives
\begin{equation}
(x_i - d_i) \cdot \kappa_j (d) = (x_j - d_j) \cdot \kappa_i (d),
\label{eq:proportion}
\end{equation}
that is, the data residuals and the consistency terms must be proportional across positions. This proportionality is the price for the existence of a single scalar: as soon as the strength ratio required at the two ends violates it, no single $\lambda$ can be optimal at both places. Boundaries requiring weak coupling and interiors requiring strong coupling are exactly the typical situation where this proportionality is broken, so the demanded profile cannot be realized by a single scalar.

The zero-residual specialization is clearer; its premise is the heterogeneity of the consistency term. If the stationary point happens to be an exact fit, namely the case where data and model agree, \eqref{eq:foc} reduces to $\lambda \cdot \kappa_i (d) = 0$ and $\lambda \cdot \kappa_j (d) = 0$. If $\lambda \neq 0$, then $\kappa_i (d) = \kappa_j (d) = 0$, which heterogeneous positions do not satisfy; if $\lambda = 0$, context simply withdraws and ill-posedness is left unattended. Neither end works, so the profile cannot be realized. The comparison with classical variational regularization theory must be stated explicitly: the linear system of a quadratic smoothness objective has a unique minimizer for every $\lambda > 0$, existence and well-posedness are unaffected, and a coupled system's solution automatically satisfies the proportionality whenever the first-order conditions hold. What this paper excludes is not the existence of solutions but the realizability of a heterogeneous demand profile by a single frozen coefficient; in the excluded configurations, a uniform $\lambda$ still gives a determinate solution, merely biased and oversmoothed relative to the ideal heterogeneous profile.\qedbox

What Lemma 1 excludes is only the most economical form; the general form stated next makes that scope exact, and whether multi-dimensional, layer-wise or position-wise parameter tables dodge the blow is left to the reach section, where dodging Lemma 1 turns out to be dodging only the shallowest negation.

The form of Lemma 1 is not what its argument uses, and the generalization a reader is entitled to ask for costs nothing. Let the data term be any additively separable term $\sum_{i \in I} f_i(d_i)$ with $f_i$ differentiable where it is evaluated, and let the consistency term be the general pairwise potential $\sum_{i \in I} \sum_{j \in I} \varphi(d_i - d_j)$ with $\varphi$ even and differentiable. Differentiating in $d_i$ and normalizing the constant factor that the evenness of $\varphi$ produces gives the stationarity condition
\begin{equation}
f_i'(d_i) + \lambda \cdot \kappa_i^{\varphi}(d) = 0, \qquad \kappa_i^{\varphi}(d) = 2 \sum_{j \in I} \varphi'(d_i - d_j),
\label{eq:focgeneral}
\end{equation}
of which \eqref{eq:foc} is the case $f_i(t) = (t - x_i)^2$, $\varphi(t) = t^2$ with $\kappa_i^{\varphi} = 4 \kappa_i$. Two features of \eqref{eq:focgeneral} decide what the generalization buys. It uses neither the squared data term nor the quadratic potential, and no convexity is used anywhere: \eqref{eq:focgeneral} is a necessary condition at a stationary point, so the exclusion below covers the stationary points of nonconvex objectives as well. And everything after the single differentiation \eqref{eq:focgeneral} is algebra in two equations, which is why the artifact separates the two steps and machine-checks both: the coordinate differentiation that produces \eqref{eq:focgeneral} from the objective on the toy carrier, and the algebra that consumes it, which is stated for a general potential.

\begin{uprop}[Lemma 1, general form]
Call $\rho_i(d) = f_i'(d_i)/\kappa_i^{\varphi}(d)$ the \emph{demand ratio} of position $i$, wherever the coefficient does not vanish. Let one and the same scalar $\lambda$ satisfy \eqref{eq:focgeneral} at positions $i$ and $j$. Then
\begin{equation}
f_i'(d_i) \cdot \kappa_j^{\varphi}(d) = f_j'(d_j) \cdot \kappa_i^{\varphi}(d),
\label{eq:proportiongeneral}
\end{equation}
which reads $\rho_i(d) = \rho_j(d)$. Consequently, if two positions have different demand ratios at an admissible configuration, then no scalar satisfies \eqref{eq:focgeneral} at both. Let now the rule be a family $(R_i)_{i \in I}$ whose values pre-exist inference, so that each $R_i = W_i$ is one number fixed before inference and independent of the instance. Then \eqref{eq:focgeneral} holds at every position at two admissible configurations $d$ and $d'$ if and only if
\begin{equation}
f_i'(d_i) \cdot \kappa_i^{\varphi}(d') = f_i'(d'_i) \cdot \kappa_i^{\varphi}(d) \qquad \text{for every position } i \in I,
\label{eq:tablecondition}
\end{equation}
that is, $\rho_i(d) = \rho_i(d')$ at every position. The minimal obstruction to a frozen family is therefore one single position at which the demand ratio changes between two admissible configurations, and the obstruction to the scalar family is the additional requirement that this ratio be the same at every position.
\end{uprop}

\argument Equation \eqref{eq:proportiongeneral} is \eqref{eq:focgeneral} at $i$ and at $j$ with $\lambda$ eliminated, the same two-substitution step that gave \eqref{eq:proportion}; no property of $f_i$ or of $\varphi$ beyond differentiability is used, which is precisely why the statement concerns stationary points rather than minimizers. For the frozen family, \eqref{eq:focgeneral} at position $i$ reads $W_i = -\rho_i(d)$ and at the same position in the second configuration $W_i = -\rho_i(d')$; the single entry $W_i$ serves both if and only if the two agree, which is \eqref{eq:tablecondition}, and the entries are chosen independently, so the condition is per position. Three remarks delimit the statement. First, the constant family is the case in which \eqref{eq:tablecondition} holds with one ratio common to all positions, so the scalar family is the shallowest object in the class rather than a separate one, and a position-wise table, however many entries it has, is not a loophole for the algebra: it satisfies \eqref{eq:focgeneral} at any single configuration by construction, and what it cannot do is follow the demand ratio across configurations, which is exactly the burden the source criterion then places on it. Second, \eqref{eq:tablecondition} is the point at which the general form meets the criterion of Section~\ref{sec:p4}: the barrier condition says that the demand separates at every scale, and the table condition is that same separation read at the granularity of the coefficient demanded at two admissible configurations, while the realization condition is what makes such a pair carry positive probability rather than consisting of a single pair of constructed instances. Third, the residue is the pair $(\varphi, f)$ itself: the argument requires the pairwise term to depend on the field through differences alone and the data term to be separable. Operator-valued couplings, in which the consistency term reads $d$ through a linear operator rather than through differences, lie outside \eqref{eq:focgeneral}, and we do not claim them.\qedbox

Two things must be separated at this point, because the indexed form of Definition~2 makes the difference visible. Lemma 1's algebra excludes only the constant family at a single configuration: a position-wise family escapes the proportionality identity \eqref{eq:proportion} there, because the stationarity conditions are then solved position by position rather than shared. Across configurations the algebra reaches further, as \eqref{eq:tablecondition} shows: one position whose demanded ratio changes between two admissible configurations excludes the frozen family however many entries it has, which is the sharpest form the algebra takes and the one the source criterion then matches. What a position-wise family does not escape is the source criterion. If its values are frozen at inference time it is Mode B by Lemma 2 and by the reach section; if its values are instantiated online from current evidence it is Mode A, and then it was never a fixed-memory rule to begin with. The index is therefore not a loophole but the setting in which both statements are stated sharply: the excluded object of the algebra is the constant family, and the excluded object of the source criterion is any family whose values pre-exist inference, constant or not.

\subsection{Lemma 2: Frozen Extrinsic Coefficients Cannot Realize a Heterogeneous Profile; the Source Is the Criterion}

Stepping back for a moment, perhaps the strength should not be hidden in a scalar but appended as an extrinsic term outside the inference objective: keep an objective faithful to the observation, and superimpose context as an additional penalty or an additional channel. This is another economical road, and with frozen coefficients it is equally impassable.

First distinguish two objective paradigms:

\begin{defn}
Call an inference objective \emph{Paradigm E} if it is an additively separable combination of a data term and an extrinsic penalty term, of the form \eqref{eq:paradigmE}; call an inference objective \emph{Paradigm I} if it is a single generative objective in which candidate structures are given by the predictions of a generative model and the objective retains only the data term itself:
\begin{equation}
\mathcal{I}(x, p) = \sum_{i \in I} (x_i - p_i)^2,
\label{eq:paradigmI}
\end{equation}
where $p$ is the prediction field given by the generative model, and no separate additive regularization term exists any longer.
\end{defn}

\begin{lemma}
Let an extrinsic rule be an extrinsic quantity whose values are frozen at inference time; in the indexed notation of Definition~2 it is a family $(R_i)_{i \in I}$ whose values pre-exist inference, and it may be constant or position-wise. Such a quantity has no effect at the level of the configuration: its optimization problem remains well-posed for any positive coefficient, with a unique solution that is merely biased and oversmoothed; if the task demands a heterogeneous strength profile, frozen coefficients cannot realize that profile at a stationary point, and if they enter the stationarity conditions they fall within the exclusion of Lemma 1, while if they do not enter they cause no causal effect in the solution. The viability of the context rule is independent of its form and depends only on the source of its values.
\end{lemma}

\argument The extrinsic side reuses the algebra of Lemma 1. The stationarity condition of Paradigm E at a fitting point requires $\lambda \cdot \kappa_i (d) = 0$ to hold simultaneously at all positions. When the consistency term is heterogeneous, that is, when there exist $i, j$ with $\kappa_i (d) \neq \kappa_j (d)$, no nonzero $\lambda$ satisfies the conditions at all positions. Treating $\lambda$ itself as a decision variable does not help either, because the system of stationarity equations itself has no solution; the extrinsic penalty leads nowhere here. The endogenous side is constructive. The objective \eqref{eq:paradigmI} of Paradigm I always has a unique minimizer: nonnegativity of the sum of squares gives $\mathcal{I}(x, p) \ge 0$ with the lower bound attained when the candidate field equals the prediction field, so minimality holds; the sum of squares vanishes if and only if every term vanishes, so uniqueness holds. Intuitively, in Paradigm I context is not a patch pasted outside the objective but a prediction grown inside the generative model; the prediction computes where context should be strong and where weak into $p$, and the data term only needs to check the prediction against the observation. To avoid a misreading, the scope of the extrinsic-side exclusion must be stated: what is excluded is the realization of a heterogeneous profile by frozen coefficients at a fitting point, not the well-posedness of the frozen extrinsic problem. A variational problem with fixed smoothing weights has a unique determinate solution, and under data-consistency constraints combined with analytic measurement models, frozen priors together with data-consistent learned reconstruction \citep{haltmeier2026} are routinely used for robust reconstruction within and, to a certain extent, beyond the training distribution. Their failure is not a lack of definition but a lack of worst-case guarantee: the solution is determinate, biased, oversmoothed, and uncontrolled on inputs beyond the corpus coverage. If an extrinsic quantity with frozen coefficients enters the stationarity conditions, it falls within the exclusion of Lemma 1; if it does not enter, it has no effect on the solution. The endogenous form, by contrast, always has a unique feasible point.\qedbox

Paradigm E and Paradigm I are two representative instances, but the form is not the criterion. The contradiction of Lemma 1 needs a single frozen scalar, that is a constant family; when position-wise weights are computed online from current evidence, the stationarity equations are solvable position by position and no contradiction arises: anisotropic diffusion and the guided filter are exactly this form \citep{perona1990,he2013}, and strengths chosen online from the current noisy observation by SURE belong to the same class \citep{stein1981}. These are Paradigm E in form and Mode A in source. What is truly excluded is frozen coefficients: whether a scalar, a parameter table, or a learned extrinsic regularizer, whenever the coefficient values are frozen offline, the extrinsic form inherits either the contradiction of Lemma 1 or the Mode B classification. For a general nonseparable but differentiable objective, the criterion likewise lands on the source rather than on the shape of the gradient field, and for a general pairwise potential with a separable data term the same conclusion follows from the general form of Lemma 1 above: whatever the potential, one position whose demanded ratio changes between two admissible configurations excludes the frozen family, and no convexity enters. The conclusion of Lemma 2 therefore rests not on an enumeration of paradigms but on a decidable property: whether the rule's values are frozen at inference time.

Lemma 2 therefore does not nail the rule's position inside a single generative objective; it converges the question onto the source: the form is free, the source is not. The only question left is ``where the rule's values come from'', which the openness criterion and the non-freezability criterion answer.

\subsection{Lemma 3: The Open Set Demands Mode A}

Inference on the open set faces a measure-theoretic problem. Let the test distribution be $P$, let the loss $\ell(R, d)$ of a rule $R$ on input $d$ measure the error of the rule's values relative to the demand rule, for example the aggregate squared error $\sum_{i \in I} \bigl(R_i(d) - t_i(d)\bigr)^2$, where $t = (t_i)_{i \in I}$ is the demand rule on that input. The demand rule deserves to be made explicit: for each input $d$ and each position $i$, $t_i(d)$ is the context strength the task actually requires at that position on that input; its measurability is secured by the integrability requirement on the loss, its existence by the content-addressing of Criterion P2, and its uniqueness in the sense of the evidence-determinacy of P2, while no identifiability assumption on $t$ is needed and Assumption A0 only requires the existence of a mapping from evidence to $t$. Let the risk be the expected loss $\mathrm{risk}(R) = \mathbb{E}_P [\ell(R, d)]$, where the expectation is taken over the distribution of inputs $d$ under $P$, let the minimum risk $\mathrm{risk}^* (R)$ be the infimum of the risks attainable by evidence-instantiated rules, and let the risk gap be
\begin{equation}
g(R) = \mathrm{risk}(R) - \mathrm{risk}^* (R).
\label{eq:gap}
\end{equation}

Generalization on the open set is taken in this paper in its strongest form: there exists a bound $\epsilon > 0$ such that $\ell(R, d) \le \epsilon$ for \emph{every} input $d$ in the test distribution. This form is stronger than boundedness in probability, but it is exactly the literal reading of ``cross-scenario reuse'': reuse means that no input makes the rule fail. The choice of standard deserves one remark: the derivation of Lemma 3 uses only the direction that the bound implies the risk does not exceed $\epsilon$. If the pointwise uniform standard is replaced by an expected-risk standard, namely $\mathrm{risk}(R) \le \epsilon$, the derivation still goes through: by \eqref{eq:gap} and nonnegativity of the minimum risk, $g(R) \le \mathrm{risk}(R) \le \epsilon$, contradicting $g(R) > \epsilon$. The negation of the decision axis therefore holds under both standards, and the choice of the pointwise strong form serves the literal reading of the narrative rather than the needs of the derivation.

\begin{lemma}
Suppose three premises hold. First, upper-tail realization: there exists an input $d_0$ such that $\ell(R, d_0) \ge \mathrm{risk}(R)$. Second, the risk gap exceeds the bound: $g(R) > \epsilon$. Third, the minimum risk is nonnegative: $\mathrm{risk}^* (R) \ge 0$. Then a generalization bound in the open-set sense cannot hold: the requirement that $\ell(R, d) \le \epsilon$ for all inputs in the test distribution is unattainable. In this sense, no rule with offline-frozen values can satisfy a uniform bound on the open set, and by the mutual exclusivity and exhaustiveness of Modes A and B, generalization requires the context rule to belong to Mode A.
\end{lemma}

\argument All three premises are listed explicitly, with no implicit assumptions. The derivation has only four lines. By the generalization requirement, $\ell(R, d_0) \le \epsilon$; by upper-tail realization, $\mathrm{risk}(R) \le \ell(R, d_0)$; together they give $\mathrm{risk}(R) \le \epsilon$. By the nonnegative minimum risk and \eqref{eq:gap}, $g(R) \le \mathrm{risk}(R)$. Combining this with $g(R) > \epsilon$ gives $\mathrm{risk}(R) > \epsilon$. The two estimates of $\mathrm{risk}(R)$ contradict each other, so the generalization requirement cannot hold.\qedbox

The conditions under which the three premises hold deserve item-by-item discussion. Upper-tail realization holds universally: for any integrable random variable, the probability that it attains a value no less than its mean is positive; this is a direct corollary of the mean inequality and depends on no distributional assumption. Nonnegativity of the minimum risk holds trivially under squared loss. What really carries content is the risk gap exceeding the bound. Criterion P4 is consumed here at the fixed tolerance of the claimed bound, and its component form is derived in Section~\ref{sec:p4}: the barrier-to-gap proposition there gives the positive floor $\kappa$ of \eqref{eq:floor} for every rule whose values pre-exist inference, so the second premise of Lemma 3 is not an independent empirical assumption but the instance of that floor at the tolerance under discussion. What remains to be tested is the satisfiability of P4 on real tasks, with two mechanistic supports. The first comes from the existing literature: the amortization gap in variational inference, namely the difference between the amortized posterior family and the per-instance optimal posterior, has been identified by theoretical and empirical studies as an attributable component of inference suboptimality \citep{cremer2018}; the fixed function family of Mode B corresponds to the amortized family, the instantiated rules of Mode A correspond to the per-instance optimum, and the two gaps are isomorphic. 

The second is a mechanism example that reduces the gap mechanism to its simplest form, and it is stated on the toy carrier of the machine verification. On that carrier the position set has four positions, position $0$ carries the evidence and position $1$ the context, so the rule family of Definition~2 has one nontrivial member, written $R$ throughout.

\begin{example}[The gap mechanism on one sample]
\label{ex:gap}
Take the three inputs $(1,0)$, $(0,1)$ and $(1,1)$. The frozen side is the one-parameter family of the toy carrier, $R(d) = W \cdot (d_0 - d_1)$, whose single coefficient exists before inference begins, so by the source clause of Definition~3 the family is Mode~B. The evidence-instantiated side consists of rules of the form $R(d) = \Phi(d_0)$, the projection of the source clause onto the function level recorded in Section~\ref{ch:axis}: the value at the position is pinned by the current evidence, and a rule implemented by an analytic expression of evidence rather than by a pre-existing statistical value is Mode~A.

\emph{A demand that reads the context.} Let the demand be $t(d) = d_0 \cdot d_1$, which varies jointly with evidence and context. The frozen family gives the values $W$, $-W$, and $0$ on the three inputs, against the demanded values $0$, $0$, and $1$; fitting the first two inputs forces $W = 0$, the third then costs the single-point squared loss $1$, and no $W$ escapes, so the total loss over the sample is at least the positive constant $1$ for every $W$ (machine-checked, \texttt{modeB\ttu positive\ttu gap}). An evidence-instantiated rule has total loss $\Phi(1)^2 + \Phi(0)^2 + (\Phi(1) - 1)^2$ on the same sample, whose minimum is $1/2$, attained at $\Phi(1) = 1/2$ and $\Phi(0) = 0$ by the constructive rule $R(d) = d_0/2$ (\texttt{modeA\ttu loss\ttu lower\ttu bound}, \texttt{modeA\ttu loss\ttu attained}), so the class gap is the positive constant $1/2$ (\texttt{toy\ttu positive\ttu gap}). The benchmark risk on this demand is $1/2 > 0$: the demand reads the context coordinate, so no evidence-instantiated rule can be exact, and the floor $\kappa$ of \eqref{eq:floor} carries a positive $\rho$ here.

\emph{A demand computable from evidence.} Keep the sample and the frozen family and replace the demand by $t(d) = d_0^2$, a function of the evidence alone, so that Assumption~A0 holds with the analytic map $\tau(x) = x^2$. The demanded values are now $1$, $0$, and $1$, while the frozen family still gives $W$, $-W$, and $0$, so its total loss is $(W - 1)^2 + W^2 + 1$, which is at least $3/2$ and is minimized at $W = 1/2$, for every $W$ (machine-checked, \texttt{modeB\ttu evidence\ttu gap}). The constructive rule $R(d) = d_0^2$ is evidence-instantiated and attains total loss $0$, hence the optimal risk (\texttt{modeA\ttu evidence\ttu attained}); the benchmark vanishes, so the whole frozen loss is gap and the gap is at least $3/2$ (\texttt{evidence\ttu gap\ttu is\ttu full}). This is the case of Assumption~A0 referred to after \eqref{eq:floor}, and there the floor reduces to $p \delta^2 / 16$ because $\rho = 0$.
\end{example}

\begin{example}[Where Criterion P4 fails]
\label{ex:p4fails}
The demand is the linear functional $t(d) = d_0 - d_1$ on the sample of Example \ref{ex:gap}. The frozen family with $W = 1$ reports exactly the demanded values, $1$, $-1$ and $0$, on the three inputs, so its total loss is $0$ and it is complete on this class at $\epsilon = 0$. The task therefore lies outside Criterion P4, and Theorem~5 is silent on it rather than contradicted by it.

What fails is the barrier condition. A demand that a fixed analytic rule pins is a demand with a pre-given modulus, here one, since $|t(d) - t(e)| \le \| d - e \|$ for every pair; the barrier condition asks the opposite, a separation $\delta$ that is fixed while the scale $\eta$ at which the instances are compared goes to zero, and for a demand of modulus one the two requirements are compatible only for $\eta \ge \delta$, so no positive $\delta$ survives the limit. The realization component is vacuous at the same time, and the risk-gap floor of \eqref{eq:floor} is not available, because the bounded-description rate $\Lambda$ is not dominated by the separation at any scale. This is the boundary of the class in its sharpest form: what fails is a premise about the task, and the linear demands referred to in Section~\ref{ch:intro} and in the limitations of Section~\ref{ch:discussion} are instances of it.
\end{example}

Two scope remarks close the example. First, the machine-checked forms are the two-sided toy of the artifact: \texttt{modeB\ttu positive\ttu gap} against \texttt{modeA\ttu loss\ttu lower\ttu bound} and \texttt{modeA\ttu loss\ttu attained} on the context-dependent demand, and \texttt{modeB\ttu evidence\ttu gap} against \texttt{modeA\ttu evidence\ttu attained} on the evidence-computable one, both decided by nonlinear arithmetic with the same axiom budget as the rest of the development (Section~\ref{ch:verify}).

Second, what the example supplies is the mechanism rather than the quantifier over rules whose values pre-exist inference, and the difference is worth stating plainly. The context-dependent demand lies outside Assumption~A0, so it corroborates that a positive gap is attainable, not the completeness of Mode A and certainly not Criterion P4 within the class; the evidence-computable demand lies inside the assumption, and there a constructive Mode A rule attains the optimal risk exactly, which is the completeness the first part cannot show. Neither part excludes every offline-frozen rule, and on a finite sample no part could: a richer family of pre-existing values fits any prescribed values, an affine form $W_0 + W_1 d_0$ matching the evidence-computable demand on these three inputs exactly. The quantifier over offline-frozen rules is supplied instead where it belongs, by the barrier condition of Section~\ref{sec:p4}: on the open set every family whose values pre-exist inference has a scale at which it cannot resolve the demand, so no such family matches every demand of a class rich in the sense of \eqref{eq:barrier}, and how often the unresolvable reversals occur is the realization level $p$ that a measurement estimates. The example is the smallest object on which the mechanism can be exhibited and machine-checked.

The two parts of the example together show that the gap mechanism can be exhibited within the class with a constructive evidence-instantiated rule attaining the optimal risk, while the satisfiability of Criterion P4 on a given task remains an empirical question; the testable statement is thereby explicit and, by Section~\ref{sec:p4}, split by component: the barrier condition is tested by comparing the demand between nearby instances across scales, the realization level is estimated from the frequency of fine-scale demand reversals, and the benchmark is the risk of the best evidence-instantiated rule available. This paper makes no magnitude claim about a concrete task: Criterion P4 is a premise about real tasks rather than a definitional clause of the class, and the numerical value of the floor \eqref{eq:floor} on a given task is an empirical measurement question left to future work.\qedbox

The direction of Lemma 3's conclusion is worth a second reading. The conditional form of the theorem does not weaken the falsification: of the three premises, two hold universally, and the risk gap exceeding the bound is carried by the component form of Criterion P4, which is derived next in this section, with twofold mechanistic support from independent literature and from the worked example of the two demands. Lemma 3 does not say that a fixed function family cannot achieve low loss on a closed benchmark; that is the weak claim under distribution agreement. What it says is that on the open set, whenever strong generalization bounded for all inputs is required, the values of the rule must be instantiated by current evidence. Combined with Lemma 2, the context rule must take its values from current evidence, while the form of the objective carrying it is unrestricted. The necessary conditions are now complete.

\subsection{The Two Components of Criterion P4}
\label{sec:p4}

Criterion P4 is the premise that carries the whole negation, and the charge that it presupposes the conclusion deserves a mechanism rather than a restatement. This subsection states P4 in two components, a barrier condition on the demand and a realization condition on the test distribution, adds a bounded-description condition that objects with pre-existing values satisfy, and derives the risk gap from the three by an elementary inequality. The purpose of the decomposition is that the conclusion of Lemma 3 is then no longer read off a criterion that asserts it. The barrier condition is a statement about the demand and mentions no rule, no average, and no benchmark; the bounded-description condition is a property of the carrying object of a rule; the realization condition is a property of the test distribution. Each of the three can fail while the other two hold, and the failure modes are the boundary cases of the task class stated below.

Fix the notation of Lemma 3. Instances $d$ range over the open set of Criterion P3, $t = (t_i )_{i \in I}$ is the demand rule, and for a profile $s = (s_i )_{i \in I}$ write $\| s \|^2 = \sum_{i \in I} s_i^2$, so that the loss of a rule on an instance is $\ell (R, d) = \| R(d) - t(d) \|^2$. The risk is $\mathrm{risk}(R) = \mathbb{E}_P [\ell (R, d)]$, the benchmark $\mathrm{risk}^*$ is the minimum risk of Lemma 3, namely the infimum of the risks attainable by evidence-instantiated rules, and the risk gap is $g(R) = \mathrm{risk}(R) - \mathrm{risk}^*$ of \eqref{eq:gap}. Let $d_{\mathrm{in}}$ be a metric on instances, and write $\mathrm{diam}_{\mathrm{in}}$ for the diameter it induces.

\paragraph{Barrier condition (P4a).} The demand has no pre-given modulus. There is a separation $\delta > 0$ such that for every scale $\eta > 0$ there are two disjoint sets of instances $A_\eta$ and $B_\eta$ with
\begin{equation}
\inf_{d \in A_\eta ,\, e \in B_\eta} \| t(d) - t(e) \| \ge \delta , \qquad \mathrm{diam}_{\mathrm{in}} (A_\eta \cup B_\eta) \le \eta .
\label{eq:barrier}
\end{equation}
The condition says that the demand separates two sets of instances that are as close together as the scale $\eta$ prescribes, with a separation that does not shrink as $\eta$ does; equivalently, the demand map is not Lipschitz on the open set with any pre-given constant. It mentions no rule: it is a property of the task alone, and it is the shape of the classical approximation barrier, in which the resolving power of a family of bounded description is compared with the oscillation of the target \citep{kolmogorov1936,pinkus1985}, and in which the massiveness of a class of functions is measured by its entropy \citep{kolmogorov1961}. Its ground on open inversion tasks is the content-addressing of Criterion P2. Where a boundary lies is carried by the content of the current evidence rather than by a pre-given grid, so a displacement smaller than any pre-given scale can carry a boundary across a position and reverse the strength demanded there. The scale $\eta$ is not bounded below by any pre-given quantity, because Criterion P3 admits inputs outside any pre-given class; a fixed-grid simulation, whose boundary positions are set by the task rather than by the content of the evidence, fails the condition, and so does any demand that is Lipschitz in the instance with a modest constant. The condition is also checkable, and cheaply: estimate the oscillation of the demand between nearby instances across resolutions and test whether it stays above a positive level.

\paragraph{Bounded-description condition (BD).} An object whose values pre-exist inference carries a finite description, and the values it can give respond to the instance at a rate bounded by that description. For a rule $R$ whose values pre-exist inference there is a finite constant $\Lambda = \Lambda (R)$, determined by the carrying parameters, such that
\begin{equation}
\| R(d) - R(e) \| \le \Lambda \cdot d_{\mathrm{in}} (d, e)
\label{eq:rate}
\end{equation}
for all instances $d$ and $e$. This is a regularity condition on implementations, and it is stated here explicitly rather than smuggled into the barrier: a feedforward computation with bounded parameters satisfies it with the layer-norm bound on its Lipschitz constant \citep{virmaux2018}, and a stored table satisfies it with the bound determined by its entries and the resolution of its index. The condition excludes one kind of object and no other: an object whose values are not carried by a finite description but computed by an analytic expression of the instance. By the implementer clause of Definition~3 such an object takes its values from an analytic invariant and belongs to Mode A, so the exclusion removes nothing from the class that the negation is about. A larger rate is not an escape either: $\Lambda$ is finite and fixed, and the barrier condition supplies a scale below which no finite rate resolves the demand.

\paragraph{Realization condition (P4b).} The separation is realized with positive frequency. There is a level $p > 0$ such that at every scale $\eta$ the sets of \eqref{eq:barrier} can be chosen so that each of them carries probability at least $p$ under the test distribution,
\begin{equation}
P (A_\eta ) \ge p , \qquad P (B_\eta ) \ge p .
\label{eq:realization}
\end{equation}
This is the empirical component of Criterion P4, and it is the component that this paper neither proves nor claims to prove: it says that the open set does not merely admit fine-scale demand reversals but presents them with a frequency bounded away from zero. It is also the component a measurement can address directly, because it is a single number attached to a task: the mass of instances on which the demand reverses across a fine scale. When it fails, the demand's fine-scale reversals are rare, and a rule whose values pre-exist inference spends most of its probability on instances where the absolute level of the demand is what a fixed description can carry; this is the honest reading of the boundary recorded for the linear-demand example of Section~\ref{ch:intro}, where a fixed-memory component is complete because the demand can be pinned down.

The three conditions are the components of Criterion P4, and the risk gap follows from them by two inequalities.

\begin{uprop}[Barrier to Gap]
Suppose the barrier condition \eqref{eq:barrier} holds with separation $\delta$, the realization condition \eqref{eq:realization} holds at level $p$, and every rule whose values pre-exist inference satisfies the bounded-description condition \eqref{eq:rate} with its own finite rate. Then every such rule satisfies
\begin{equation}
\mathrm{risk}(R) \ge \frac{p \delta^2}{16} , \qquad \text{and hence} \qquad g(R) \ge \frac{p \delta^2}{16} - \rho = \kappa ,
\label{eq:floor}
\end{equation}
where $\rho$ is any upper bound on the benchmark risk $\mathrm{risk}^*$. In particular, if $\rho < p \delta^2 / 16$, then $\kappa > 0$, the risk gap is at least $\kappa$ for every rule whose values pre-exist inference, and no generalization bound in the sense of Lemma 3 holds at any tolerance $\epsilon < \kappa$.
\end{uprop}

\argument Fix a rule $R$ whose values pre-exist inference and let $\Lambda$ be its rate. Choose a scale $\eta \le \delta / (2 \Lambda)$ and the sets $A = A_\eta$, $B = B_\eta$ of \eqref{eq:barrier}. For $d \in A$ and $e \in B$, the triangle inequality gives
\begin{equation}
\| R(d) - t(d) \| + \| R(e) - t(e) \| \ge \| t(d) - t(e) \| - \| R(d) - R(e) \| \ge \delta - \Lambda \eta \ge \frac{\delta}{2},
\label{eq:deficit}
\end{equation}
where the middle step uses \eqref{eq:barrier} and \eqref{eq:rate}. Hence at least one member of every pair has a deficit of at least $\delta / 4$. If some $d \in A$ and some $e \in B$ both had a deficit below $\delta / 4$, that inequality would fail for the pair; therefore the set $S$ of instances whose deficit is at least $\delta / 4$ contains all of $A$ or all of $B$. By \eqref{eq:realization} it carries probability at least $p$, so $\mathrm{risk}(R) \ge p \, (\delta / 4)^2$. The bound on $g(R)$ follows from \eqref{eq:gap} and $\mathrm{risk}^* \le \rho$, and the final sentence repeats the derivation of Lemma 3 with $\epsilon < \kappa$.\qedbox

Five remarks delimit what the decomposition does and does not buy.

First, the gap is now derived rather than assumed, and what carries the derivation is a statement about the demand alone. The barrier condition is about the demand and can fail on a task that satisfies P1 to P3, the linear-demand example of Section~\ref{ch:intro} being such a task; the bounded-description condition is about implementations; the realization condition is about the test distribution. A criterion that quantified over all tolerances in a single assertion is thereby a weak structural condition, a regularity condition, and one measurable number.

Second, the decomposition is not in conflict with universal approximation theorems, and the point deserves a sentence because the objection is natural. Those theorems approximate a fixed class of targets by a sequence of descriptions of growing size; \eqref{eq:floor} concerns one description of pre-given size, uniformly over the open set and at every scale, and asserts nothing against a sequence. What carries the force is the finiteness of the description together with the unbounded resolution of the barrier: for every pre-given rate there is a scale at which the demand outruns it. That is exactly the situation of a fixed-memory component, whose values exist before inference begins and whose description is therefore finite and fixed.

Third, what the decomposition yields is a positive floor rather than an unbounded gap, and the criterion has to be read in its component form rather than as a single quantified assertion. A risk gap exceeding \emph{every} pre-given tolerance is stronger than the barrier supports and, for a bounded loss, stronger than is satisfiable; what \eqref{eq:floor} gives is the constant $\kappa$ determined by the task's separation $\delta$, its realization level $p$, and the benchmark $\rho$. No force is lost in the argument, because the generalization claim that Lemma 3 negates is made at a fixed tolerance: what is falsified is a uniform bound at every tolerance below $\kappa$. Under Assumption A0 the benchmark risk vanishes in the limit of exact evidence instantiation, and $\kappa$ then reduces to $p \delta^2 / 16$; the context-dependent demand of the first part of Example~\ref{ex:gap} is deliberately not of that kind, which is why its benchmark is a positive constant rather than zero, while the evidence-computable demand of the second part is.

Fourth, the honest sentence that the force of the theorems equals the force of the premise is retained in a sharper form. The force of Theorem 5 equals exactly the force of Criterion P4; what has changed is that P4 is now a sufficient condition implying a bound, with its three components named separately and each addressable by a different means: the barrier by comparing the demand across scales, the realization level by measuring the frequency of the reversals, and the benchmark by measuring the risk of the best evidence-instantiated rule available. The floor is also the precise quantitative content of the word ``completes'' in the claim that is falsified: since the demand separates by $\delta$ at every scale, a rule whose risk gap is at least $\kappa$ misses the task by at least the separation the task imposes on itself, so a bound at a tolerance at or above $\kappa$ is not a completion of the task but a statement weaker than the task's own structure requires.

Fifth, the dichotomy is not what the floor rests on, and the relaxation most often proposed does not create a third mode. An implementation that is Mode~A except for a finite budget of frozen metaparameters --- a learned step size, a fixed dictionary, a corpus-fitted basis --- is still Mode~B by the source criterion, and it is already inside the quantitative statement rather than outside it: it carries a finite description, so the bounded-description condition holds for it at some finite rate $\Lambda$, and the barrier supplies the scale $\eta \le \delta / (2\Lambda)$ at which its values cannot follow the demand, whence the floor $\kappa$ of \eqref{eq:floor} applies to it as to any other frozen family. The bookkeeping exemption of Definition~3 is not in tension with this, because it exempts quantities that do not change the rule's values, whereas the frozen budget under discussion is what determines them and is therefore classified. What carries the floor is the finiteness of the description together with the barrier, not the axis itself, so no hybrid escapes by being hybrid.

\subsection{Classification and Closure: Neural Networks}

What remains is only classification. Inference by a neural network feeds the current input into a set of weights in a forward computation; the weights do not change at inference time, and their values come from the statistics of a training corpus or from manual setting. The two sources do not change the classification:

\begin{proposition}
The context rule of a neural network, insofar as its values vary with corpus statistics or manual setting, is carried by weights fixed at inference time and belongs to Mode B.
\end{proposition}

\argument The network's rule satisfies \eqref{eq:modeB}, with the parameter object $W$ being the weights and the computational processes $F_i$ being forward propagation. The weights already exist when inference begins, and whether their values were produced by gradient descent or by manual setting does not change their prior existence at inference time. Membership in Mode B follows from the definition. The more substantive part is the converse: a rule carried by fixed weights does not belong to Mode A. Take the collision sample pair given by Criterion P3, namely two inputs $d$ and $e$ with identical evidence and freely bifurcating context, satisfying $d|_{\ev} = e|_{\ev}$ while $d$ and $e$ take different values at context positions. If the rule belonged to Mode A, then by \eqref{eq:modeA} it would have to give the same value at every position on the two inputs, because the evidence is the same; but a nondegenerate rule with fixed parameters gives different values on the two inputs, because the value varies with the context part. The two requirements conflict, so the rule does not belong to Mode A. What the collision argument excludes are rules that are nondegenerate with respect to evidence-context bifurcation: a network that computes only an analytic function of the evidence is not within its reach, but by the implementer clause of Definition~3, the source of the rule in that case is an analytic invariant rather than the component's memory, and the completeness is not credited to the component.\qedbox

The nondegeneracy clause is exact: the degenerate rule with $W = 0$ is excluded by the hypothesis, and a rule with $W$ arbitrarily close to zero cannot be convicted by the collision argument alone; by the component form of Criterion P4 in Section~\ref{sec:p4}, such rules still fall within the reach of the negation, merely through a different argument path.

Mode A and Mode B are mutually exclusive, the necessary condition given by Lemma 3 is violated, and closure is immediate:

\begin{theorem}
Neural networks alone cannot independently complete open inversion tasks; Claim $\mathcal{P}$ is false.
\end{theorem}

\argument Open inversion tasks satisfy Criterion P3, so collision sample pairs exist; Proposition 4 gives that the network's rule belongs to Mode B and not to Mode A; if Claim $\mathcal{P}$ were true, there would be a bound $\epsilon$ under which the network rule satisfies the uniform bound, and the component form of Criterion P4 in Section~\ref{sec:p4} gives a risk gap bounded below by the positive floor $\kappa$ of \eqref{eq:floor} for every rule whose values pre-exist inference, so at every tolerance $\epsilon < \kappa$ the second premise of Lemma 3 is satisfied and Lemma 3 in turn gives that generalization requires the rule to belong to Mode A. The necessary condition conflicts with the nature of the carrier, and no uniform bound holds at any tolerance below the floor of \eqref{eq:floor}; Claim $\mathcal{P}$ is false.\qedbox

\subsection{Approximate Inference: Cost and Accuracy, Not Provenance}

Approximate inference is the natural reply to the provenance question and deserves a direct answer, because every modern inference pipeline performs it. Loopy belief propagation replaces the exact marginal with a message-passing fixed point on the graph \citep{pearl1988,murphy1999,yedidia2005}. Structured variational inference replaces it with an optimum over a tractable family of distributions on a substructure \citep{saul1996,wainwright2008}. A learned proposal replaces the hand-designed message or family with a function fitted to a corpus \citep{mnih2014}. The three differ in cost and in accuracy, and none of them touches the source of the rule's values.

The point is formal rather than diagrammatic. Replace the exact evaluation of the rule by an approximate one: the rule is still a family in the sense of Definition~2 whose values are given offline, and the operator's own quantities --- a message schedule, a damping factor, an iteration budget --- are bookkeeping in the sense of Definition~3, because they affect the faithfulness with which the rule is evaluated and not the rule itself, while a learned proposal carries corpus-fitted values of its own and its rule is Mode~B for the same reason. Hence \eqref{eq:modeB} continues to hold, the source criterion places the approximate rule in Mode~B exactly as before, and \eqref{eq:floor} applies to it with its own bounded-description rate: approximation cannot restore the uniform bound of Lemma~3, because whatever accuracy it buys, the values remain frozen before inference and the floor is a property of the source rather than of the operator. The accuracy of the approximation enters the risk in addition to the source-side error, with a sign that is not controlled: loopy belief propagation is not guaranteed to converge on graphs with cycles and can possess multiple fixed points, and a variational family fixes its own gap to the per-instance optimum, which is the amortization gap of Section~\ref{ch:related} \citep{cremer2018} and an instance of the Mode~B gap rather than a way around it. Approximation therefore relocates the failure from the cost of evaluating the rule to the accuracy of the evaluation; it does not move the rule across the decision axis. If the proposal or the family is instead instantiated from current evidence --- an analytic proposal, a message computed from the current measurement, a family parameter read off the evidence --- then a Mode~A mechanism has entered and the credit belongs to it rather than to the component. As with the fresh randomness of diffusion sampling in the reach section below, the operator's internal quantities carry no information about the demand and do not constitute a third source of values.

\subsection{Reach}

Finally, delimit the reach, stating what the negation covers and what it does not.

Lemma 1 excludes only the single global scalar form. Networks employing multi-dimensional, layer-wise, or position-wise smoothing parameters seem to remain in the gap, but in fact they do not. Once the parameter field is frozen, its values are still given by the corpus, and it still belongs to Mode B, merely a Mode B of higher dimension; and this is not only the source criterion speaking, because the algebra of the general form \eqref{eq:tablecondition} already excludes the frozen family as soon as one position changes its demanded ratio between two admissible configurations, whatever the number of entries. The classification of Theorem 5 is insensitive to the shape of the parameter object: one scalar is a parameter table, and a million parameters are also a parameter table; the verdict rests on whether the table's values pre-exist and are fixed at inference time, not on the size of the table. Heterogeneity is content-addressed: where boundaries appear depends on the current scene, and a frozen parameter table can only record where boundaries \emph{usually} are, not where they are \emph{in this frame}. Dodging Lemma 1 is dodging only the shallowest negation.

The opposite of a frozen parameter table deserves to be stated positively: implementations whose position-wise weights are computed online from current evidence, such as anisotropic diffusion, the guided filter \citep{perona1990,he2013}, and strengths chosen online by SURE \citep{stein1981}, are neither frozen parameter tables nor subject to the contradiction of Lemma 1, because the stationarity equations are solvable position by position. They are Paradigm E in form and Mode A in source, and they lie outside the negation. The only legitimate way past Lemma 1 is that values are given online by evidence.

One might let the parameter field itself be generated from evidence by another network; hypernetworks are exactly this idea \citep{ha2016}. The generator's weights are still frozen: if their values come from a corpus, Mode B is merely moved up one level, the rule's values are now given by a ``value generator of frozen values'', and tracing back any number of levels, the ultimate source of the values is still the corpus, with recursion creating no new source; if the generator itself is an analytic function, for example sampling weights computed by projective geometry, then the rule's values are instantiated by evidence while the source is an analytic invariant, and by the implementer clause of Definition~3 the rule belongs to Mode A and lies outside the negation. Test-time adaptation is analogous: the parameters updated online are given by a frozen update rule, so a corpus-learned update rule merely moves Mode B up one level, while an analytic closed-form update rule, such as a Kalman gain, belongs to Mode A. Recursion creates no new source; it only lengthens the chain of tracing. Likewise, attaching classical extrinsic smoothing priors to learnable modules \citep{lan2023}, or plugging conditional random fields into networks as recurrent layers \citep{zheng2015}, does not change the source of the rule's values; the difference in craftsmanship between appending and internalizing does not touch the decision axis.

There is a subtler escape hatch worth naming. A frozen-weight network that computes only an analytic function of the evidence, such as a geometric calculation substituting camera intrinsics, depends only on evidence at the function level, and the collision argument cannot touch it. By the implementer clause of Definition~3, the rule of such an implementation is still classified by the source of the function: if the function is an analytic invariant it belongs to Mode A, and if the function is learned from a corpus it belongs to Mode B. The escape hatch therefore does not overturn the decision axis but corroborates it: a network implementing only an analytic invariant has not completed the task as a fixed-memory component; it is merely the carrier of the invariant, and the credit for completeness does not go to the component's memory.

Two families of mainstream contemporary implementations must be named. Plug-and-play priors \citep{venkatakrishnan2013} and diffusion posterior sampling \citep{chung2023} combine a frozen prior with likelihood guidance driven online by the current measurement in one system: the frozen prior belongs to Mode B and carries no guarantee, while guidance strengths and step sizes, if instantiated by current evidence, credit completeness to the Mode A part. Data-consistent learned reconstruction takes this combination one step further into a systematic form: frozen components propose candidates, measurement consistency is enforced by analytic constraints, and determinate solutions are obtained that are robust within and to a certain extent beyond the training distribution \citep{haltmeier2026}. This corroborates the credit decomposition of this chapter: data consistency supplies a verifiable analytic criterion, and the frozen component supplies candidates without a guarantee. The fresh randomness injected by diffusion sampling at inference time is neither a corpus value nor an evidence value; it carries no information about the demand, so it does not constitute a third source of values and does not change the classification.

In-context learning is of the same type: a frozen-weight model treats the current context as evidence at inference time, and its output is a function of the evidence at the function level; by the implementer clause, a function learned from a corpus belongs to Mode B and earns no exemption from its behaviorally online adaptation.

What this paper falsifies has never been a network of some shape, but the carrier itself that bears rules with parameters fixed at inference time. The only implementations genuinely left outside the negation are those whose values at inference time come only from current evidence or analytic invariants, which belong to Mode A to begin with. Nor is what this paper falsifies ``these components are useless'', but only ``solo use is complete''.
